\textbf{Phần 1. Trắc nghiệm lựa chọn}
\setcounterpageref{ex}{0}

\begin{ex} %[3X1V80]
Bảng thống kê số bạn thích các môn thể thao: Bóng đá: 25; Bơi lội: 18; Cầu lông: 22; Cờ vua: 15. Môn thể thao được nhiều bạn thích nhất là
\choice
{Bơi lội}
{Cầu lông}
{Cờ vua}
{\True Bóng đá}
\loigiai{
Bóng đá có 25 bạn thích, nhiều nhất.
\textbf{Đáp án của câu hỏi là D}
}
\end{ex}

\begin{ex} %[3X1V80]
Tổng số bạn tham gia khảo sát ở câu trên là
\choice
{$75$ bạn}
{\True $80$ bạn}
{$85$ bạn}
{$70$ bạn}
\loigiai{
Tổng: $25 + 18 + 22 + 15 = 80$ bạn.
\textbf{Đáp án của câu hỏi là B}
}
\end{ex}

\begin{ex} %[3X1V80]
Số bạn thích Bóng đá nhiều hơn số bạn thích Cờ vua bao nhiêu bạn?
\choice
{$8$ bạn}
{$12$ bạn}
{\True $10$ bạn}
{$7$ bạn}
\loigiai{
$25 - 15 = 10$ bạn.
\textbf{Đáp án của câu hỏi là C}
}
\end{ex}

\begin{ex} %[3X1V80]
Trong hộp chỉ có 6 chiếc bút màu xanh. Lấy ra ngẫu nhiên 1 chiếc bút. Sự kiện nào là 'Chắc chắn' xảy ra?
\choice
{\True Lấy được bút màu xanh}
{Lấy được bút màu đỏ}
{Lấy được bút màu đen}
{Không lấy được bút nào}
\loigiai{
Trong hộp chỉ có bút xanh nên chắc chắn lấy được bút xanh.
\textbf{Đáp án của câu hỏi là A}
}
\end{ex}

\begin{ex} %[3X1V80]
Gieo 1 con xúc xắc chuẩn. Sự kiện nào là 'Không thể' xảy ra?
\choice
{Xuất hiện mặt 1 chấm}
{Xuất hiện mặt 6 chấm}
{Xuất hiện mặt 3 chấm}
{\True Xuất hiện mặt 0 chấm}
\loigiai{
Xúc xắc chỉ có từ 1 đến 6 chấm nên không thể có mặt 0 chấm.
\textbf{Đáp án của câu hỏi là D}
}
\end{ex}

\begin{ex} %[3X1V80]
Trong túi có 4 kẹo cam và 4 kẹo dâu. Lấy ra 1 kẹo. Khả năng lấy được kẹo cam là
\choice
{Chắc chắn xảy ra}
{\True Có thể xảy ra}
{Không thể xảy ra}
{Luôn luôn xảy ra}
\loigiai{
Có cả 2 loại kẹo nên có thể lấy được kẹo cam.
\textbf{Đáp án của câu hỏi là B}
}
\end{ex}

\begin{ex} %[3X1V80]
Thống kê điểm kiểm tra Toán của tổ 1: có 3 điểm 10, 4 điểm 9, 2 điểm 8, 1 điểm 7. Số học sinh của tổ 1 là
\choice
{$9$ bạn}
{$11$ bạn}
{\True $10$ bạn}
{$8$ bạn}
\loigiai{
Tổng số bạn: $3 + 4 + 2 + 1 = 10$ bạn.
\textbf{Đáp án của câu hỏi là C}
}
\end{ex}

\begin{ex} %[3X1V80]
Số điểm từ 9 trở lên trong tổ 1 là bao nhiêu điểm?
\choice
{\True $7$ điểm}
{$6$ điểm}
{$8$ điểm}
{$9$ điểm}
\loigiai{
Điểm 9 và 10 gồm: $3 + 4 = 7$ điểm.
\textbf{Đáp án của câu hỏi là A}
}
\end{ex}

\begin{ex} %[3X1V80]
Một hộp có 10 viên bi xanh và 2 viên bi đỏ. Cần lấy ra ít nhất bao nhiêu viên bi để CHẮC CHẮN có ít nhất 1 viên bi màu đỏ?
\choice
{$3$ viên bi}
{$10$ viên bi}
{$12$ viên bi}
{\True $11$ viên bi}
\loigiai{
Trường hợp xấu nhất lấy cả 10 bi xanh, cần lấy thêm 1 viên nữa: $10 + 1 = 11$ viên.
\textbf{Đáp án của câu hỏi là D}
}
\end{ex}

\begin{ex} %[3X1V80]
Ngày mai có tuyết rơi giữa mùa hè tại Hà Nội là sự kiện
\choice
{Chắc chắn xảy ra}
{\True Không thể xảy ra (hoặc cực kì khó xảy ra)}
{Thường xuyên xảy ra}
{Luôn xảy ra}
\loigiai{
Khí hậu Hà Nội mùa hè nắng nóng, không thể có tuyết rơi.
\textbf{Đáp án của câu hỏi là B}
}
\end{ex}

\begin{ex} %[3X1V80]
Bảng theo dõi số sách đọc được trong hè: An: 12 cuốn; Bình: 15 cuốn; Chi: 18 cuốn. Trung bình mỗi bạn đọc được bao nhiêu cuốn sách?
\choice
{$14$ cuốn}
{$16$ cuốn}
{\True $15$ cuốn}
{$18$ cuốn}
\loigiai{
$(12 + 15 + 18) : 3 = 45 : 3 = 15$ cuốn.
\textbf{Đáp án của câu hỏi là C}
}
\end{ex}

\begin{ex} %[3X1V80]
Tung một đồng xu 30 lần, mặt Sấp xuất hiện 14 lần. Mặt Ngửa xuất hiện bao nhiêu lần?
\choice
{\True $16$ lần}
{$14$ lần}
{$15$ lần}
{$17$ lần}
\loigiai{
$30 - 14 = 16$ lần.
\textbf{Đáp án của câu hỏi là A}
}
\end{ex}

\textbf{Phần 2. Trắc nghiệm đúng sai}
\setcounterpageref{ex}{0}

\begin{ex} %[3X1V80]
Xét tính đúng sai của việc phân tích bảng số liệu và xác suất:
\choiceTF
{\True Bảng số liệu giúp tổng hợp thông tin rõ ràng.}
{Một sự kiện chắc chắn xảy ra thì có lúc không xảy ra.}
{\True Trong trò chơi gieo xúc xắc, mặt xuất hiện số chẵn là sự kiện có thể xảy ra.}
{Lấy 1 viên bi từ hộp chỉ có bi đỏ thì có thể lấy được bi xanh.}
\loigiai{
\begin{itemize}
    \item Ý a) Đúng: Lợi ích bảng thống kê.
    \item Ý b) Sai: Sai vì chắc chắn xảy ra thì luôn luôn xảy ra.
    \item Ý c) Đúng: Có các mặt 2, 4, 6.
    \item Ý d) Sai: Sai vì không thể lấy được bi xanh.
\end{itemize}
\begin{itemize}
    \item a) \textbf{Đúng}.
    \item b) \textbf{Sai}.
    \item c) \textbf{Đúng}.
    \item d) \textbf{Sai}.
\end{itemize}
}
\end{ex}

\begin{ex} %[3X1V80]
Bảng thống kê số cây trồng được của 3 lớp: 3A trồng 45 cây, 3B trồng 55 cây, 3C trồng 50 cây:
\choiceTF
{\True Lớp 3B trồng được nhiều cây nhất.}
{\True Lớp 3A trồng được ít cây nhất.}
{\True Cả 3 lớp trồng được tất cả 150 cây.}
{Lớp 3B trồng gấp đôi lớp 3A.}
\loigiai{
\begin{itemize}
    \item Ý a) Đúng: 55 cây nhiều nhất.
    \item Ý b) Đúng: 45 cây ít nhất.
    \item Ý c) Đúng: $45 + 55 + 50 = 150$ cây.
    \item Ý d) Sai: Sai vì $45 \times 2 = 90 \neq 55$.
\end{itemize}
\begin{itemize}
    \item a) \textbf{Đúng}.
    \item b) \textbf{Đúng}.
    \item c) \textbf{Đúng}.
    \item d) \textbf{Sai}.
\end{itemize}
}
\end{ex}

\textbf{Phần 3. Trả lời ngắn}
\setcounterpageref{ex}{0}

\begin{ex} %[3X1V80]
Một hộp có 8 bi xanh và 4 bi đỏ. Cần lấy ra ít nhất bao nhiêu viên bi để chắc chắn có ít nhất 1 viên bi màu đỏ?
\shortans{9}
\loigiai{
\textbf{Đáp số:} $9$\\
Trường hợp xấu nhất lấy 8 bi xanh, lấy thêm 1: $8 + 1 = 9$ viên.
}
\end{ex}

\begin{ex} %[3X1V80]
Tổng số học sinh tham gia 3 câu lạc bộ là 90 bạn: vẽ có 30 bạn, võ có 35 bạn. Hỏi câu lạc bộ nhạc có bao nhiêu bạn?
\shortans{25}
\loigiai{
\textbf{Đáp số:} $25$\\
$90 - (30 + 35) = 25$ bạn.
}
\end{ex}

\begin{ex} %[3X1V80]
Gieo đồng xu 40 lần, mặt ngửa xuất hiện 18 lần. Mặt sấp xuất hiện bao nhiêu lần?
\shortans{22}
\loigiai{
\textbf{Đáp số:} $22$\\
$40 - 18 = 22$ lần.
}
\end{ex}

\begin{ex} %[3X1V80]
Số điểm 10 của tổ 1 là 28 điểm, tổ 2 là 32 điểm. Cả hai tổ có tất cả bao nhiêu điểm 10?
\shortans{60}
\loigiai{
\textbf{Đáp số:} $60$\\
$28 + 32 = 60$ điểm.
}
\end{ex}

\begin{ex} %[3X1V80]
Có bao nhiêu mặt của con xúc xắc chuẩn ghi số chấm bé hơn 4?
\shortans{3}
\loigiai{
\textbf{Đáp số:} $3$\\
Mặt 1, 2, 3 (có 3 mặt).
}
\end{ex}

\begin{ex} %[3X1V80]
Số cây trồng của 4 bạn: An: 5 cây; Bình: 7 cây; Cúc: 6 cây; Dũng: 6 cây. Tính tổng số cây của cả 4 bạn.
\shortans{24}
\loigiai{
\textbf{Đáp số:} $24$\\
$5 + 7 + 6 + 6 = 24$ cây.
}
\end{ex}
